\chapter{Independence and Zero--One Laws}
\begin{center}\itshape
The harmonic series $1+\frac12+\frac13+\cdots$ diverges, and $1-\frac12+\frac13-\cdots$ converges to $\log2$.\\
Choose every sign by tossing a fair coin. With probability one the series converges,\\
and the density of its sum at $2$ is within $10^{-42}$ of $\frac18$.
\end{center}

In 1909 \'Emile Borel published a paper with a provocative title, \emph{Les probabilit\'es d\'enombrables et leurs applications
arithm\'etiques}. Until then probability had dealt with finitely many trials, or with continuous quantities treated by
geometry. Borel asked about infinitely many: what is the probability that, in an infinite sequence of coin tosses, heads
eventually stops appearing? That the frequency of heads tends to one half? He answered, and deduced that almost every real
number is \emph{normal}, every digit appearing with its fair frequency in every base. His arguments had gaps, and he knew
it: nobody had yet said what ``the probability of an event involving infinitely many tosses'' means.

Fourteen years later Hugo Steinhaus gave the answer. A sequence of fair coin tosses is the same thing as the binary digits
of a point chosen uniformly in $[0,1]$, and the probability of an event is the Lebesgue measure of the set of points
whose digits realise it. The year before, Hans Rademacher had proved, in the language of orthogonal functions, that the
series $\sum a_n\varepsilon_n$ with signs $\varepsilon_n$ given by the binary digits converges for almost every point when
$\sum a_n^2<\infty$; Khintchine and Kolmogorov proved the converse in 1925. And in 1933 Andrei Kolmogorov's
\emph{Grundbegriffe der Wahrscheinlichkeitsrechnung} put the whole subject on measure theory: a probability is a measure
of total mass one, an event is a measurable set, independence is a product formula. In an appendix of that small book he
proved a strange theorem. An event that depends on a sequence of independent variables, but on no finite part of it, has
probability $0$ or $1$. Nothing in between.

The zero--one law is a theorem about knowledge: it tells you that a question has a definite answer before telling you
what the answer is. Does a random series converge? Is the radius of convergence of a random power series larger than
$\frac12$? Does a random walk come back to $0$ infinitely often? Each of these has probability $0$ or $1$, and the work
is to find out which. In 1955 Edwin Hewitt and Leonard Savage extended the law to events that may depend on the first
variables, provided they are unchanged by any shuffle of finitely many of them, and this is exactly what a random walk
needs. This chapter proves both laws, shows that the convergence of a series of independent terms in probability already
forces almost sure convergence (Paul L\'evy), and settles the random sign problem of Rademacher, Khintchine and
Kolmogorov. It ends with the random harmonic series of the epigraph and its smooth, strange density.

\begin{reminder}[from Tomes 3 and 6]
\emph{Independence} (Tome~3, Chapter~5; Tome~6, Chapter~6). Events $A_1,\dots,A_n$ are independent if
$P(\bigcap_{i\in I}A_i)=\prod_{i\in I}P(A_i)$ for every $I\subseteq\{1,\dots,n\}$. \emph{Dynkin's $\pi$--$\lambda$
theorem} (Tome~6, Theorem~6.4): a $\lambda$-system (closed under complements and countable disjoint unions, containing
$\Omega$) that contains a $\pi$-system (closed under finite intersections) contains the $\sigma$-algebra it generates.
\emph{Existence} (Tome~6, Theorems~6.17 and~11.8): for any probability measures $\mu_1,\mu_2,\dots$ on $\R$ there is a
probability space carrying independent random variables $X_n$ of laws $\mu_n$. \emph{Borel--Cantelli} (Tome~3,
Theorem~4.11): if $\sum P(A_n)<\infty$ then $P(A_n\text{ i.o.})=0$; if the $A_n$ are independent and $\sum P(A_n)=\infty$
then $P(A_n\text{ i.o.})=1$. \emph{Martingales} (Tome~3, Chapter~10): Kolmogorov's inequality $P(\max_{k\le n}|S_k|\ge
a)\le\operatorname{Var}(S_n)/a^2$ (Corollary~10.16) and the convergence of martingales bounded in $L^2$ (Theorem~10.18).
\emph{Characteristic functions} (Tome~3, Chapter~6): convergence in distribution implies pointwise convergence of
characteristic functions, and an integrable characteristic function $\varphi$ has the continuous density
$\frac1{2\pi}\int e^{-itx}\varphi(t)\dd t$ (Theorem~6.8).
\end{reminder}

\section{Independent $\sigma$-algebras}

Throughout, $(\Omega,\mathcal F,P)$ is a probability space. Sub-$\sigma$-algebras $\mathcal G_1,\dots,\mathcal G_n$ of
$\mathcal F$ are \emph{independent} if $P(G_1\cap\dots\cap G_n)=P(G_1)\cdots P(G_n)$ whenever $G_i\in\mathcal G_i$; an
infinite family is independent if each finite subfamily is. Random variables are independent when the $\sigma$-algebras
$\sigma(X_i)=\{X_i^{-1}(B)\}$ they generate are. Checking the product formula on all of $\mathcal G_i$ is rarely possible;
the following lemma says that generators are enough.

\begin{lemma}[$\pi$-systems suffice]\label{lem:pisuffice}
Let $\mathcal P_1,\dots,\mathcal P_n$ be $\pi$-systems with $\Omega\in\mathcal P_i$, such that $P(A_1\cap\dots\cap A_n)=
\prod P(A_i)$ for all $A_i\in\mathcal P_i$. Then $\sigma(\mathcal P_1),\dots,\sigma(\mathcal P_n)$ are independent.
\end{lemma}
\begin{proof}
Fix $A_2\in\mathcal P_2,\dots,A_n\in\mathcal P_n$ and let $\mathcal L$ be the class of $A\in\mathcal F$ with
$P(A\cap A_2\cap\dots\cap A_n)=P(A)P(A_2)\cdots P(A_n)$. It contains $\Omega$ and $\mathcal P_1$. It is closed under
complements: subtract the identity for $A$ from the identity for $\Omega$. It is closed under countable disjoint unions, by
countable additivity on both sides. So $\mathcal L$ is a $\lambda$-system containing $\mathcal P_1$, hence $\sigma(
\mathcal P_1)$, by Dynkin's theorem. Thus the hypothesis holds with $\mathcal P_1$ replaced by $\sigma(\mathcal P_1)$, which is
again a $\pi$-system containing $\Omega$. Repeat with the second index, and so on. Taking some $A_i=\Omega$ gives the
product formula for every subfamily.
\end{proof}

For a random variable, $\{X\le x\}$, $x\in\R$, together with $\Omega$, is a $\pi$-system generating $\sigma(X)$, so
$X_1,\dots,X_n$ are independent as soon as $P(X_1\le x_1,\dots,X_n\le x_n)=\prod P(X_i\le x_i)$: the definition of Tome~3
agrees with this one.

\begin{proposition}[grouping]\label{prop:grouping}
Let $(\mathcal G_i)_{i\in I}$ be independent, and let $I=\bigsqcup_jI_j$ be a partition. Then the $\sigma$-algebras
$\mathcal H_j=\sigma\bigl(\bigcup_{i\in I_j}\mathcal G_i\bigr)$ are independent.
\end{proposition}
\begin{proof}
Let $\mathcal P_j$ be the class of finite intersections $G_{i_1}\cap\dots\cap G_{i_r}$ with distinct $i_k\in I_j$ and
$G_{i_k}\in\mathcal G_{i_k}$. It is a $\pi$-system containing $\Omega$ and generating $\mathcal H_j$. For finitely many
distinct $j$, a product of elements of the $\mathcal P_j$ is a finite intersection of sets from distinct $\mathcal G_i$, so
its probability is the product of the probabilities, which regroups as $\prod_jP(\text{element of }\mathcal P_j)$. Lemma~\ref{lem:pisuffice}
 concludes.
\end{proof}

So if $X_1,X_2,\dots$ are independent, then $f(X_1,\dots,X_5)$, $g(X_6,X_7)$ and $h(X_8,X_9,\dots)$ are independent for all
measurable $f,g,h$, a fact used constantly and silently from now on.

\section{Kolmogorov's zero--one law}

\begin{definition}
Let $X_1,X_2,\dots$ be random variables. The \emph{tail $\sigma$-algebra} is
\[
\mathcal T=\bigcap_{n\ge1}\sigma(X_n,X_{n+1},\dots).
\]
An event of $\mathcal T$ is a \emph{tail event}: for each $n$, whether it happens is decided by $X_n,X_{n+1},\dots$ alone.
\end{definition}

With $S_n=X_1+\dots+X_n$, the event $\{\sum X_n\text{ converges}\}$ is a tail event, since convergence does not depend on
the first terms. So are $\{X_n>1\text{ i.o.}\}$ and $\{\limsup S_n/n\le c\}$, since $S_n/n$ and $(S_n-S_m)/n$ have the same
$\limsup$. The event $\{S_n>0\text{ i.o.}\}$ is not a tail event in general: changing $X_1$ can change it.

\begin{theorem}[Kolmogorov's zero--one law, 1933]\label{thm:kolmogorov01}
If $X_1,X_2,\dots$ are independent, every tail event has probability $0$ or $1$.
\end{theorem}
\begin{proof}
Let $A\in\mathcal T$ and $n\ge1$. Then $A$ lies in $\sigma(X_{n+1},X_{n+2},\dots)$, which is independent of $\sigma(X_1,\dots,X_n)$ by
grouping. So $A$ is independent of every event of the union $\mathcal A$ of the $\sigma$-algebras $\sigma(X_1,\dots,X_n)$, which is a
$\pi$-system (an increasing union of $\sigma$-algebras) and generates $\sigma(X_1,X_2,\dots)$, so by Lemma~\ref{lem:pisuffice},
$A$ is independent of $\sigma(X_1,X_2,\dots)$, which contains $A$. Hence $P(A)=P(A\cap A)=P(A)^2$.
\end{proof}

The proof is a little scandalous: $A$ is independent of itself. A tail event is so far in the future that knowing all
of the present tells you nothing about it, and yet it is determined by the whole sequence.

\begin{corollary}\label{cor:tailconst}
If $X_1,X_2,\dots$ are independent and $Y$ is a $\mathcal T$-measurable random variable with values in $[-\infty,\infty]$,
then $Y$ is almost surely constant. In particular, $\limsup S_n/a_n$ and $\liminf S_n/a_n$ are almost surely constant for
every sequence $a_n\to\infty$, and $\sum X_n$ converges with probability $0$ or $1$.
\end{corollary}
\begin{proof}
The function $c\mapsto P(Y\le c)$ takes only the values $0$ and $1$, is non-decreasing and right-continuous; let
$c_0=\inf\{c:P(Y\le c)=1\}\in[-\infty,\infty]$. Then $P(Y\le c)=1$ for $c>c_0$ and $0$ for $c<c_0$, so $Y=c_0$ almost surely.
For the examples, $\limsup_n S_n/a_n=\limsup_n(S_n-S_m)/a_n$ for each $m$, which is $\sigma(X_{m+1},\dots)$-measurable.
\end{proof}

\begin{think}
Let $X_n$ be independent with $P(X_n=n)=P(X_n=-n)=\frac12$. Is $\{\lim S_n/n^2\text{ exists}\}$ a tail event? What is its
probability? Now let $X_n$ be independent and uniform on $[0,1]$: what is the probability that infinitely many of them are
below $\frac1n$? Below $\frac1{n^2}$? The zero--one law guarantees the answers are $0$ or $1$; Borel--Cantelli finds them.
\end{think}

\section{Exchangeable events and random walks}

A \emph{finite permutation} is a bijection $\pi$ of $\{1,2,\dots\}$ that moves only finitely many integers. It acts on
sequences by $(\pi x)_i=x_{\pi(i)}$. An event $A\in\sigma(X_1,X_2,\dots)$ is \emph{exchangeable} if it is unchanged by every
finite permutation of the variables: $A=\{(X_n)\in B\}$ for a measurable $B\subseteq\R^{\N}$ with $\pi^{-1}B=B$ for every
finite $\pi$. Tail events are exchangeable, but so is $\{S_n>0\text{ i.o.}\}$: a permutation of $X_1,\dots,X_m$ leaves
$S_n$ unchanged for $n\ge m$.

We need one approximation lemma.

\begin{lemma}\label{lem:approx}
Let $\mathcal A$ be an algebra of sets (closed under complements and finite unions) and $A\in\sigma(\mathcal A)$. For every
$\varepsilon>0$ there is $A'\in\mathcal A$ with $P(A\,\triangle\,A')<\varepsilon$.
\end{lemma}
\begin{proof}
Let $\mathcal L$ be the class of sets that can be so approximated. It contains $\mathcal A$, and is closed under complements
since $A^c\triangle A'^c=A\triangle A'$. If $A_1,A_2,\dots\in\mathcal L$ are disjoint with union $A$, choose $N$ with
$\sum_{k>N}P(A_k)<\varepsilon/2$ and $A'_k\in\mathcal A$ with $P(A_k\triangle A_k')<\varepsilon/2N$; then
$A'=A_1'\cup\dots\cup A_N'\in\mathcal A$ and $A\triangle A'\subseteq\bigcup_{k\le N}(A_k\triangle A_k')\cup\bigcup_{k>N}A_k$,
of probability less than $\varepsilon$. So $\mathcal L$ is a $\lambda$-system containing the $\pi$-system $\mathcal A$, and
Dynkin's theorem gives $\mathcal L\supseteq\sigma(\mathcal A)$.
\end{proof}

\begin{theorem}[Hewitt--Savage, 1955]\label{thm:hewittsavage}
If $X_1,X_2,\dots$ are independent and identically distributed, every exchangeable event has probability $0$ or $1$.
\end{theorem}
\begin{proof}
Let $A=\{X\in B\}$ be exchangeable, where $X=(X_n)$, and let $\varepsilon>0$. By Lemma~\ref{lem:approx} applied to the
algebra $\bigcup_n\sigma(X_1,\dots,X_n)$, there are $n$ and a measurable $C\subseteq\R^n$ such that $A_n=\{(X_1,\dots,X_n)\in
C\}$ satisfies $P(A\triangle A_n)<\varepsilon$. Let $\pi$ be the finite permutation exchanging $k$ and $n+k$ for $1\le k\le n$,
and let $B_n=\{(X_{n+1},\dots,X_{2n})\in C\}=\{\pi X\in C\times\R^{\N}\}$.

The sequences $X$ and $\pi X$ have the same law, since they are both i.i.d.\ with the same marginal (the law of an i.i.d.\
sequence is determined by its finite-dimensional distributions, Tome~6, Theorem~11.8). Since $\pi^{-1}B=B$, the event
$A$ is also $\{\pi X\in B\}$, so
\[
P(A\triangle B_n)=P\bigl(\pi X\in B\triangle(C\times\R^\N)\bigr)=P\bigl(X\in B\triangle(C\times\R^\N)\bigr)=P(A\triangle A_n)<\varepsilon.
\]
Now $A_n$ and $B_n$ are independent, depending on disjoint blocks. Since $|P(E)-P(F)|\le P(E\triangle F)$, and $(A_n\cap
B_n)\triangle A\subseteq(A_n\triangle A)\cup(B_n\triangle A)$,
\[
|P(A)-P(A)^2|\le|P(A)-P(A_n\cap B_n)|+|P(A_n)P(B_n)-P(A)^2|\le2\varepsilon+2\varepsilon.
\]
As $\varepsilon$ is arbitrary, $P(A)=P(A)^2$.
\end{proof}

The zero--one law is about \emph{events}; its most famous consequence is about \emph{paths}.

\begin{theorem}[the fate of a random walk]\label{thm:trichotomy}
Let $X_1,X_2,\dots$ be i.i.d.\ and $S_n=X_1+\dots+X_n$. Exactly one of the following holds almost surely:
\begin{enumerate}
\item $S_n=0$ for all $n$ (this is the case $X_1=0$ a.s.);
\item $S_n\to+\infty$;
\item $S_n\to-\infty$;
\item $\liminf S_n=-\infty$ and $\limsup S_n=+\infty$.
\end{enumerate}
\end{theorem}
\begin{proof}
The random variable $L=\limsup S_n$ is measurable for the exchangeable events: $\{L\le c\}$ is unchanged by a finite
permutation, which changes only finitely many $S_n$. By Hewitt--Savage, $L=c$ almost surely for a constant
$c\in[-\infty,\infty]$. Let $S'_n=X_2+\dots+X_{n+1}$. The sequence $(S'_n)$ has the same law as $(S_n)$, so $\limsup S'_n=c$
almost surely too. But $S_{n+1}=X_1+S'_n$, so $c=X_1+c$ almost surely. If $c$ is finite, then $X_1=0$ almost surely, which
is case~1. Otherwise $c=\pm\infty$; similarly $\liminf S_n=c'\in\{\pm\infty\}$. The combinations $c=-\infty$ (so $S_n\to-\infty$),
$c'=+\infty$ (so $S_n\to+\infty$) and $c'=-\infty<c=+\infty$ are cases 3, 2 and~4.
\end{proof}

\begin{figure}[ht]
\centering\fig{ch01-trichotomy}
\caption{Three random walks with Gaussian steps of variance $1$ and means $0.12$, $0$ and $-0.12$. Theorem~\ref{thm:trichotomy} says
these are the only possible fates. With a mean, the strong law decides; without one, the walk swings ever wider in both
directions (Corollary~\ref{cor:oscillate}).}
\label{fig:trichotomy}
\end{figure}

Which case occurs? If $\mathbb E|X_1|<\infty$ and $\mathbb E X_1>0$, the strong law of large numbers (Tome~3, Corollary~4.12 in
the bounded case, Chapter~2 in general) gives $S_n/n\to\mathbb E X_1>0$, hence case~2. Centred walks oscillate:

\begin{corollary}\label{cor:oscillate}
If $\mathbb E X_1=0$ and $0<\mathbb E X_1^2<\infty$, then $\limsup S_n/\sqrt n=+\infty$ and $\liminf S_n/\sqrt n=-\infty$
almost surely. In particular, a centred walk with finite variance crosses every level infinitely often.
\end{corollary}
\begin{proof}
Let $\sigma^2=\mathbb E X_1^2$ and $M>0$. The event $E=\{S_n>M\sqrt n\text{ i.o.}\}$ is exchangeable. By Fatou's lemma applied to
indicators, $P(E)\ge\limsup_nP(S_n>M\sqrt n)$, and by the central limit theorem this is $P(Z>M/\sigma)>0$ for a standard
Gaussian $Z$. So $P(E)>0$, hence $P(E)=1$. Taking $M=1,2,3,\dots$ gives $\limsup S_n/\sqrt n=+\infty$; the $\liminf$ is the
same with $-X$.
\end{proof}

How fast does the walk swing? Not like $\sqrt n$, which it exceeds by every factor infinitely often, but only a little more:
the law of the iterated logarithm of Chapter~3 says $\limsup S_n/\sqrt{2n\log\log n}=\sigma$.

\section{Series of independent terms}

For independent summands, three modes of convergence of $\sum X_n$ coincide: almost sure, in probability, and (Chapter~4)
in distribution. The first equivalence rests on a maximal inequality.

\begin{lemma}[Ottaviani's inequality]\label{lem:ottaviani}
Let $X_1,\dots,X_n$ be independent, $S_k=X_1+\dots+X_k$, $a>0$, and $c=\max_{k\le n}P(|S_n-S_k|\ge a)<1$. Then
\[
P\Bigl(\max_{k\le n}|S_k|\ge2a\Bigr)\le\frac{P(|S_n|\ge a)}{1-c}.
\]
\end{lemma}
\begin{proof}
Let $\tau$ be the first $k\le n$ with $|S_k|\ge2a$ ($\tau=\infty$ if there is none). The event $\{\tau=k\}$ depends on
$X_1,\dots,X_k$, and $S_n-S_k$ on $X_{k+1},\dots,X_n$. On $\{\tau=k\}\cap\{|S_n-S_k|<a\}$ we have $|S_n|\ge|S_k|-|S_n-S_k|>a$.
These events are disjoint for different $k$, so
\[
\begin{gathered}
P(|S_n|\ge a)\ge\sum_{k\le n}P(\tau=k)\,P(|S_n-S_k|<a)\\\ge(1-c)\sum_{k\le n}P(\tau=k)=(1-c)\,P\Bigl(\max_{k\le n}|S_k|\ge2a\Bigr).
\end{gathered}
\]
Divide by $1-c$.
\end{proof}

\begin{theorem}[L\'evy, 1937]\label{thm:levyequiv}
Let $X_1,X_2,\dots$ be independent. If $S_n$ converges in probability, it converges almost surely.
\end{theorem}
\begin{proof}
Let $S_n\to S$ in probability, and let $\varepsilon,\delta\in(0,\frac12)$. Since $|S_j-S_i|\le|S_j-S|+|S-S_i|$, there is $m$
such that $P(|S_j-S_i|\ge\varepsilon)\le\delta$ for all $j\ge i\ge m$. Apply Ottaviani's inequality to the independent variables
$X_{m+1},\dots,X_n$, whose partial sums are $S_k-S_m$: here $c\le\delta$, so
\[
P\Bigl(\max_{m\le k\le n}|S_k-S_m|\ge2\varepsilon\Bigr)\le\frac{\delta}{1-\delta}\le2\delta.
\]
Let $n\to\infty$: $P(\sup_{k\ge m}|S_k-S_m|\ge2\varepsilon)\le2\delta$, hence $P(\sup_{k,l\ge m}|S_k-S_l|\ge4\varepsilon)\le2\delta$. The
event $D_\varepsilon=\bigcap_m\{\sup_{k,l\ge m}|S_k-S_l|\ge4\varepsilon\}$ therefore has probability at most $2\delta$ for every
$\delta$, so $P(D_\varepsilon)=0$. The sequence $(S_n)$ fails to be Cauchy exactly on $\bigcup_{j\ge1}D_{1/j}$, a null set.
\end{proof}

Independence is essential: Exercise~\ref{ex:levyindep} gives dependent summands whose partial sums converge in probability but
nowhere. With Tome~3 we already have the simplest sufficient condition:

\begin{proposition}[Kolmogorov's $L^2$ criterion]\label{prop:l2series}
If $X_n$ are independent, $\mathbb E X_n=0$ and $\sum\mathbb E X_n^2<\infty$, then $\sum X_n$ converges almost surely and in $L^2$.
\end{proposition}
\begin{proof}
$S_n$ is a martingale for the filtration $\sigma(X_1,\dots,X_n)$, and $\mathbb E S_n^2=\sum_{k\le n}\mathbb E X_k^2$ is bounded,
so Tome~3, Theorem~10.18 applies.
\end{proof}

The next chapter turns this into a necessary and sufficient condition for all independent series, Kolmogorov's three-series
theorem. Here we settle the most symmetric case first.

\section{Random signs}

Let $\varepsilon_1,\varepsilon_2,\dots$ be independent with $P(\varepsilon_n=1)=P(\varepsilon_n=-1)=\frac12$: \emph{Rademacher
variables}, the binary digits of a uniform point in disguise. For real $a_n$, when does $\sum a_n\varepsilon_n$ converge?
Proposition~\ref{prop:l2series} says: if $\sum a_n^2<\infty$. The converse needs a lower bound on the size of a sum of random
signs, which we get from a simple and widely useful inequality.

\begin{lemma}[Paley--Zygmund]\label{lem:pz}
If $Y\ge0$ and $0<\mathbb E Y^2<\infty$, then for $0\le\theta<1$,
\[
P(Y>\theta\,\mathbb E Y)\ge(1-\theta)^2\frac{(\mathbb E Y)^2}{\mathbb E Y^2}.
\]
\end{lemma}
\begin{proof}
$\mathbb E Y=\mathbb E[Y\mathbf 1_{Y\le\theta\mathbb E Y}]+\mathbb E[Y\mathbf 1_{Y>\theta\mathbb E Y}]\le\theta\,\mathbb E Y+
\sqrt{\mathbb E Y^2\,P(Y>\theta\mathbb E Y)}$ by Cauchy--Schwarz. Rearrange and square.
\end{proof}

\begin{lemma}[moments of a Rademacher sum]\label{lem:radmoments}
For real $a_1,\dots,a_n$ and $T=\sum_{k\le n}a_k\varepsilon_k$, $\sigma^2=\sum a_k^2$: $\mathbb E T^2=\sigma^2$, $\mathbb E T^4\le3\sigma^4$,
and $P(T^2>\frac12\sigma^2)\ge\frac1{12}$.
\end{lemma}
\begin{proof}
Expand: $\mathbb E[\varepsilon_i\varepsilon_j\varepsilon_k\varepsilon_l]$ is $1$ when the indices pair up and $0$ otherwise. So
$\mathbb E T^4=\sum_ia_i^4+3\sum_{i\ne j}a_i^2a_j^2\le3\bigl(\sum a_i^2\bigr)^2$. Apply Paley--Zygmund to $Y=T^2$ with $\theta=\frac12$:
$P(T^2>\frac12\sigma^2)\ge\frac14\cdot\frac{\sigma^4}{3\sigma^4}$.
\end{proof}

\begin{theorem}[Rademacher, 1922; Khintchine--Kolmogorov, 1925]\label{thm:rademacher}
The series $\sum a_n\varepsilon_n$ converges almost surely if $\sum a_n^2<\infty$, and diverges almost surely if $\sum a_n^2=\infty$.
\end{theorem}
\begin{proof}
The first half is Proposition~\ref{prop:l2series}. Suppose $\sum a_n^2=\infty$, and suppose the series converges with positive
probability. Convergence is a tail event, so the series converges almost surely, and its partial sums $S_n$ form a Cauchy
sequence in probability: there is $m_0$ with $P(|S_n-S_m|\ge\frac12)<\frac1{12}$ for all $n>m\ge m_0$. But since $\sum a_k^2$
diverges, there is $n>m_0$ with $\sigma^2=\sum_{m_0<k\le n}a_k^2\ge1$, and Lemma~\ref{lem:radmoments} applied to
$T=S_n-S_{m_0}$ gives $P(|T|>\frac1{\sqrt2})\ge\frac1{12}$, a contradiction since $\frac1{\sqrt2}>\frac12$. So the series
converges with probability $0$.
\end{proof}

So $\sum\pm\frac1{n^\alpha}$ with random signs converges almost surely exactly when $\alpha>\frac12$, against $\alpha>1$ for
absolute convergence and $\alpha>0$ for alternating signs. Random signs sit in between: their cancellation is the square
root cancellation of the central limit theorem, and the dividing line $\sum a_n^2$ is the variance.

\begin{remark}[random power series]
Theorem~\ref{thm:rademacher} applied to $\sum a_n\varepsilon_nz^n$ at each fixed $z$, together with the zero--one law, shows that
the radius of convergence of a random power series is almost surely constant (Exercise~\ref{ex:radius} does the i.i.d.\ case).
Much more is true: if $\sum a_nz^n$ has radius $1$, then almost surely $\sum\varepsilon_na_nz^n$ cannot be continued analytically
across any arc of the unit circle (a theorem of Paley and Zygmund; Steinhaus had proved the analogue for random phases). We shall not need this; see
Kahane's book in the Sources, Chapter~4.
\end{remark}

\section{The random harmonic series}

Let $X=\sum_{n\ge1}\frac{\varepsilon_n}n$. By Theorem~\ref{thm:rademacher}, it converges almost surely, and $\mathbb E X=0$,
$\mathbb E X^2=\sum\frac1{n^2}=\frac{\pi^2}6$. What does its law look like?

\begin{proposition}\label{prop:harmonic}
The characteristic function of $X$ is $\varphi(t)=\prod_{n\ge1}\cos\frac tn$, and for $|t|\ge\pi$,
\[
|\varphi(t)|\le e^{-\pi|t|/6}.
\]
Consequently $X$ has a density $g$ which is even, of class $C^\infty$, and given by $g(x)=\frac1\pi\int_0^\infty\cos(xt)\varphi(t)\dd t$.
\end{proposition}
\begin{proof}
$\mathbb E e^{it\varepsilon_n/n}=\cos\frac tn$, so the partial sum $S_N$ has characteristic function $\prod_{n\le N}\cos\frac tn$.
Since $S_N\to X$ almost surely, $e^{itS_N}\to e^{itX}$, and dominated convergence gives the product formula.

For the bound we use $\cos x\le e^{-x^2/2}$ for $|x|\le\frac\pi2$: the function $h(x)=\log\cos x+\frac{x^2}2$ has $h(0)=0$ and
$h'(x)=x-\tan x\le0$ on $[0,\frac\pi2)$. Let $t\ge\pi$ (the function is even) and $M=\lfloor2t/\pi\rfloor$. For $n>M$ we have
$t/n<\frac\pi2$, and for $n\le M$ we bound $|\cos\frac tn|\le1$. So
\[
|\varphi(t)|\le\exp\Bigl(-\frac{t^2}2\sum_{n>M}\frac1{n^2}\Bigr)\le\exp\Bigl(-\frac{t^2}2\cdot\frac1{M+1}\Bigr)\le
\exp\Bigl(-\frac{t^2}{2(2t/\pi+1)}\Bigr)\le e^{-\pi t/6},
\]
using $\sum_{n>M}n^{-2}\ge\int_{M+1}^\infty x^{-2}\dd x$ and, in the last step, $2t/\pi+1\le3t/\pi$ for $t\ge\pi$.

So $|t|^k\varphi(t)$ is integrable for every $k$. By Tome~3, Theorem~6.8, $X$ has the continuous density
$g(x)=\frac1{2\pi}\int e^{-ixt}\varphi(t)\dd t$, which equals $\frac1\pi\int_0^\infty\cos(xt)\varphi(t)\dd t$ because $\varphi$ is even
and real. Differentiating $k$ times under the integral is allowed, since $|t^k\varphi(t)|$ is an integrable bound independent of $x$
(Tome~6, Theorem~7.17). Finally $g$ is even because $X$ and $-X$ have the same law.
\end{proof}

\begin{figure}[ht]
\centering\fig{ch01-harmonic}
\caption{The density $g$ of $\sum\pm\frac1n$, computed from Proposition~\ref{prop:harmonic}. It is flat-topped near $0$, where
$g(0)=0.249994\ldots$, just below $\frac14$, and at $x=2$ it passes through $\frac18$ to within $10^{-42}$ (Schmuland, 2003),
without being exactly $\frac18$. Direct computation gives both values.}
\label{fig:harmonic}
\end{figure}

\begin{think}
The density in Figure~\ref{fig:harmonic} is almost constant on $[-\frac12,\frac12]$, where it stays between $0.2495$ and $0.25$. Look at the first two terms:
$\varepsilon_1+\frac{\varepsilon_2}2$ takes the four values $\pm\frac12,\pm\frac32$. Can you guess why the density of the sum
should be nearly flat on $[-\frac12,\frac12]$, and why $g(x)$ near $x=2$ should be close to half of its value near $0$? (Schmuland's
paper investigates the near miss at $x=2$; it is a beautiful piece of numerical detective work.)
\end{think}

\begin{beauty}
Zero--one laws give certainty for free. Without computing anything, they tell you that a random series converges almost surely or
diverges almost surely, that a random walk drifts away or swings forever, that a random power series has a deterministic radius.
The price is that they never tell you which. The rest of this book is, in a sense, the art of paying that price: of finding, for each
question the zero--one laws declare decided, the probability that decides it.
\end{beauty}

\section{Exercises}

\begin{exercise}[$\star$]\label{ex:tailevents}
Let $X_1,X_2,\dots$ be independent and $S_n=X_1+\dots+X_n$. Which of the following events are tail events? (a) $\{\sum X_n$ converges$\}$;
(b) $\{\sum X_n$ converges and its sum is $0\}$; (c) $\{X_n>1$ i.o.$\}$; (d) $\{S_n>0$ i.o.$\}$; (e) $\{\lim S_n/n$ exists and equals $5\}$.
For those that are not, find independent $X_n$ for which the event has probability strictly between $0$ and~$1$.
\end{exercise}

\begin{exercise}[$\star$]\label{ex:subgauss}
Let $X=\sum_n\varepsilon_n/n$. Prove that $\mathbb E e^{\lambda X}=\prod_n\cosh\frac\lambda n\le e^{\pi^2\lambda^2/12}$ and deduce
$P(X\ge x)\le e^{-3x^2/\pi^2}$ for $x\ge0$. How does this compare with $P(X\ge4)$ for a Gaussian of the same variance?
\end{exercise}

\begin{exercise}[$\star$]\label{ex:polya}
Hewitt--Savage needs independence. In P\'olya's urn (Tome~3, Theorem~10.21), start with one red and one green ball; at each step draw
a ball, return it with another of the same colour, and let $X_n=1$ if the $n$-th ball drawn is red. Show that the sequence $(X_n)$ is
exchangeable (its law is invariant under finite permutations), but that the exchangeable event $\{\lim\frac1n\sum_{k\le n}X_k>\frac12\}$
has probability $\frac12$.
\end{exercise}

\begin{exercise}[$\star\star$]\label{ex:levyindep}
L\'evy's theorem needs independence. On $[0,1]$ with Lebesgue measure, list the dyadic intervals $[0,1],[0,\frac12],[\frac12,1],
[0,\frac14],\dots$ as $I_1,I_2,\dots$, and let $S_n=\mathbf 1_{I_n}$, $X_1=S_1$, $X_n=S_n-S_{n-1}$. Show that $S_n\to0$ in probability but
that $S_n(\omega)$ converges for no $\omega$.
\end{exercise}

\begin{exercise}[$\star\star$]\label{ex:radius}
Let $X_0,X_1,\dots$ be i.i.d., not almost surely $0$, and let $R$ be the radius of convergence of $\sum X_nz^n$. Show that $R$ is almost surely
constant, that $R=1$ almost surely if $\mathbb E\log^+|X_0|<\infty$, and that $R=0$ almost surely otherwise.
\end{exercise}

\begin{exercise}[$\star\star$]\label{ex:infmean}
Let $X_1,X_2,\dots$ be i.i.d.\ with $\mathbb E|X_1|=\infty$. Show that $\limsup|S_n|/n=\infty$ almost surely. (So the strong law of large
numbers cannot hold with any finite limit: the integrability assumption of Chapter~2 is necessary.)
\end{exercise}

\begin{exercise}[$\star\star\star$]\label{ex:unbounded}
Let $\sum a_n^2=\infty$ and $S_n=\sum_{k\le n}a_k\varepsilon_k$. Theorem~\ref{thm:rademacher} says $S_n$ diverges almost surely. Prove the
stronger statement $\limsup S_n=+\infty$ and $\liminf S_n=-\infty$ almost surely, even when $a_n\to0$.
\end{exercise}

\section{Hints}

\begin{hint}{ex:tailevents}
Changing one $X_k$ changes the sum of a convergent series. For (d), take $X_1$ with two values and $X_n=0$ for $n\ge2$.
\end{hint}
\begin{hint}{ex:subgauss}
Compare the Taylor series of $\cosh y$ and $e^{y^2/2}$ term by term, then use Markov's inequality on $e^{\lambda X}$ and optimise in $\lambda$.
\end{hint}
\begin{hint}{ex:polya}
Compute the probability of a given sequence of colours: it depends only on the numbers of reds and greens. Then use the uniform limit of
Tome~3.
\end{hint}
\begin{hint}{ex:levyindep}
Every $\omega$ lies in infinitely many of the $I_n$ and outside infinitely many.
\end{hint}
\begin{hint}{ex:radius}
$1/R=\limsup|X_n|^{1/n}$. For $c>1$, $\sum P(|X_0|>c^n)$ is finite or infinite according as $\mathbb E\log^+|X_0|$ is.
\end{hint}
\begin{hint}{ex:infmean}
$\sum_nP(|X_1|>Cn)\ge\mathbb E|X_1|/C-1$, and $X_n=S_n-S_{n-1}$.
\end{hint}
\begin{hint}{ex:unbounded}
Use Lemma~\ref{lem:radmoments} on $S_n$ itself to show that $\{\sup_n|S_n|=\infty\}$ has positive probability, then the zero--one law, then
the symmetry $\varepsilon\mapsto-\varepsilon$.
\end{hint}

\section{Solutions}

\begin{solution}{ex:tailevents}
\idea A tail event is one that cannot be changed by changing finitely many variables, while staying a measurable function of the rest.
Test each event by modifying $X_1$.
\writeup (a) is a tail event: for each $m$, convergence of $\sum_nX_n$ is equivalent to convergence of $\sum_{n>m}X_n$. (c) is a tail
event: $\{X_n>1\text{ i.o.}\}=\{X_n>1\text{ for infinitely many }n>m\}$ for every $m$. (e) is a tail event, because $S_n/n-(S_n-S_m)/n=
S_m/n\to0$, so the event is $\{\lim(S_n-S_m)/n=5\}$, measurable with respect to $\sigma(X_{m+1},\dots)$.

(b) is not: take $X_1=\pm1$ with probability $\frac12$ each, $X_2=1$ and $X_n=0$ for $n\ge3$; then $\sum X_n=X_1+1$ vanishes with
probability $\frac12$. (d) is not: take $X_1=\pm1$ with probability $\frac12$ and $X_n=0$ for $n\ge2$; then $S_n=X_1$ for all $n$ and
$P(S_n>0\text{ i.o.})=\frac12$. (Both are exchangeable events, but these variables are not identically distributed.)
\end{solution}

\begin{solution}{ex:subgauss}
\idea Bound the moment generating function by that of a Gaussian, then use the exponential Markov inequality, as in Chernoff's method
(Tome~3, Theorem~4.6).
\writeup $\cosh y=\sum\frac{y^{2k}}{(2k)!}\le\sum\frac{y^{2k}}{2^kk!}=e^{y^2/2}$, since $(2k)!\ge2^kk!$. By independence,
$\mathbb E e^{\lambda S_N}=\prod_{n\le N}\cosh\frac\lambda n\le\exp\bigl(\frac{\lambda^2}2\sum_n\frac1{n^2}\bigr)=e^{\pi^2\lambda^2/12}$. Since
$e^{\lambda S_N}\to e^{\lambda X}$ almost surely, Fatou's lemma gives $\mathbb E e^{\lambda X}\le\liminf_N\mathbb E e^{\lambda S_N}\le e^{\pi^2\lambda^2/12}$.
(The equality $\mathbb E e^{\lambda X}=\prod\cosh\frac\lambda n$ also holds, by uniform integrability, but the inequality suffices.) For $\lambda>0$,
Markov's inequality gives $P(X\ge x)\le e^{-\lambda x}\mathbb E e^{\lambda X}\le e^{-\lambda x+\pi^2\lambda^2/12}$, which is smallest at $\lambda=6x/\pi^2$,
where it equals $e^{-3x^2/\pi^2}$.

For $x=4$ the bound is $e^{-48/\pi^2}\approx7.7\times10^{-3}$. A Gaussian with variance $\pi^2/6$, standard deviation $1.2825$, exceeds $4$ with
probability $P(N(0,1)\ge3.119)\approx9.1\times10^{-4}$. The truth is much smaller still: integrating the density of Proposition~\ref{prop:harmonic},
one finds numerically $P(X\ge4)\approx2.0\times10^{-7}$. The tails of $X$ are much thinner than Gaussian, and this exercise captures only their
Gaussian part: the variance $\frac{\pi^2}6$ comes mostly from the first few terms, and reaching $4$ needs many terms to conspire.
\end{solution}

\begin{solution}{ex:polya}
\idea Exchangeability is about the law of finite strings of colours; Hewitt--Savage fails because the colours are not independent, they
share an unknown limiting proportion.
\writeup After $r$ reds and $g$ greens have been drawn, the urn holds $r+1$ red and $g+1$ green balls among $r+g+2$. Hence the probability of any
particular sequence of $n$ draws with $r$ reds and $g=n-r$ greens is
\[
\frac{1\cdot2\cdots r\cdot1\cdot2\cdots g}{2\cdot3\cdots(n+1)}=\frac{r!\,g!}{(n+1)!},
\]
since the red numerators are $1,2,\dots,r$ in whatever order they occur, and likewise the green ones. This depends only on $r$, so the law
of $(X_1,\dots,X_n)$ is invariant under permutations; this holds for every $n$, so the law of the sequence is invariant under finite
permutations. The proportion of red balls is a martingale converging almost surely to a uniform random variable $U$ on $[0,1]$ (Tome~3,
Theorem~10.21), and $\frac1n\sum_{k\le n}X_k$ has the same limit (the number of reds drawn is the number of reds in the urn minus one). So
the event has probability $P(U>\frac12)=\frac12$. It is exchangeable, since the limit of averages is unchanged by permuting finitely many
terms.
\end{solution}

\begin{solution}{ex:levyindep}
\idea The typewriter sequence moves an indicator across $[0,1]$ again and again with shrinking width: small in measure, but every point is
visited forever.
\writeup The intervals of generation $j$ have length $2^{-j}$ and cover $[0,1]$. So $P(S_n\ne0)=|I_n|\to0$, which is convergence in
probability to $0$. For each $\omega$ and each generation $j$, some interval of generation $j$ contains $\omega$, and some (for $j\ge1$)
does not, so $S_n(\omega)=1$ and $S_n(\omega)=0$ both happen infinitely often. The $X_n$ are of course dependent.
\end{solution}

\begin{solution}{ex:radius}
\idea The radius is given by the Cauchy--Hadamard formula, a $\limsup$ of the terms, so it is a tail variable; then the question is how often
$|X_n|$ exceeds $c^n$, a Borel--Cantelli question.
\writeup $1/R=\limsup|X_n|^{1/n}$ is unchanged by changing finitely many $X_n$, so it is tail measurable and almost surely constant by
Corollary~\ref{cor:tailconst}. For $c>0$, $\{|X_n|^{1/n}>c\}=\{\log|X_n|>n\log c\}$. For $Y\ge0$ we have $\sum_{n\ge1}P(Y>n)\le\mathbb E Y\le
\sum_{n\ge0}P(Y>n)$, by taking expectations in $\sum_{n\ge1}\mathbf 1_{Y>n}\le Y\le\sum_{n\ge0}\mathbf 1_{Y>n}$. Take $c>1$ and $Y=\log^+|X_0|/\log c$: if $\mathbb E\log^+|X_0|<\infty$ then $\sum P(|X_n|>c^n)<\infty$, so by
Borel--Cantelli $|X_n|^{1/n}\le c$ eventually, and $\limsup|X_n|^{1/n}\le1$. Conversely, $P(X_0\ne0)>0$ gives $\delta>0$ with $P(|X_0|>\delta)>0$,
so $|X_n|>\delta$ infinitely often (second Borel--Cantelli), and $\delta^{1/n}\to1$ gives $\limsup|X_n|^{1/n}\ge1$. So $R=1$. If
$\mathbb E\log^+|X_0|=\infty$, then for every $c>1$, $\sum P(|X_n|>c^n)=\infty$, and by the second Borel--Cantelli lemma $|X_n|^{1/n}>c$
infinitely often; so $\limsup|X_n|^{1/n}=\infty$ and $R=0$.
\end{solution}

\begin{solution}{ex:infmean}
\idea If the terms themselves are often larger than $Cn$, the partial sums cannot all stay below $Cn/2$.
\writeup Fix $C>0$. Since $\mathbb E|X_1|/C=\mathbb E[|X_1|/C]\le\sum_{n\ge0}P(|X_1|>Cn)$, the series $\sum_nP(|X_n|>Cn)=\sum_nP(|X_1|>Cn)$ diverges.
The events $\{|X_n|>Cn\}$ are independent, so by the second Borel--Cantelli lemma $|X_n|>Cn$ infinitely often. Since $|X_n|\le|S_n|+|S_{n-1}|$,
for each such $n$ either $|S_n|>Cn/2$ or $|S_{n-1}|>Cn/2\ge C(n-1)/2$. So $\limsup|S_n|/n\ge C/2$ almost surely, for every $C$.
\end{solution}

\begin{solution}{ex:unbounded}
\idea Divergence alone allows bounded oscillation. To get unboundedness, show directly that $|S_n|$ is large with probability bounded below
(the Paley--Zygmund bound does exactly this), turn ``with positive probability'' into ``almost surely'' with a zero--one law, and use the symmetry
of the signs to split between $+\infty$ and $-\infty$.
\writeup Let $\sigma_n^2=\sum_{k\le n}a_k^2\to\infty$. By Lemma~\ref{lem:radmoments}, $P(|S_n|>\sigma_n/\sqrt2)\ge\frac1{12}$ for every $n$. For any
$M$, $\{|S_n|>M\}\supseteq\{|S_n|>\sigma_n/\sqrt2\}$ once $\sigma_n\ge\sqrt2M$, so $P(|S_n|>M\text{ i.o.})\ge\limsup_nP(|S_n|>M)\ge\frac1{12}$ (Fatou
for indicators). Hence $E=\{\sup_n|S_n|=\infty\}$ has probability at least $\frac1{12}$. Changing finitely many signs changes every $S_n$ by at most a
fixed amount, so $E$ is a tail event, and $P(E)=1$.

Now $E\subseteq\{\limsup S_n=+\infty\}\cup\{\liminf S_n=-\infty\}$. Both events are tail events by the same remark, and they have the same probability
$p$, since $(\varepsilon_n)$ and $(-\varepsilon_n)$ have the same law. So $1\le2p$, which forces $p=1$ by the zero--one law.
\end{solution}

\section*{Sources and further reading}

\begin{itemize}
\item \emph{History.} \'E.~Borel, ``Les probabilit\'es d\'enombrables et leurs applications arithm\'etiques'', \emph{Rend.\ Circ.\ Mat.\ Palermo} 27
(1909), 247--271. H.~Steinhaus, ``Les probabilit\'es d\'enombrables et leur rapport \`a la th\'eorie de la mesure'', \emph{Fund.\ Math.}\ 4 (1923),
286--310. H.~Rademacher, ``Einige S\"atze \"uber Reihen von allgemeinen Orthogonalfunktionen'', \emph{Math.\ Ann.}\ 87 (1922), 112--138.
A.~Khintchine and A.~Kolmogorov, ``\"Uber Konvergenz von Reihen, deren Glieder durch den Zufall bestimmt werden'', \emph{Mat.\ Sb.}\ 32 (1925), 668--677.
A.~N.~Kolmogorov, \emph{Grundbegriffe der Wahrscheinlichkeitsrechnung}, Springer, 1933 (the zero--one law is in the appendix); English translation
\emph{Foundations of the Theory of Probability}, Chelsea, 1956. E.~Hewitt and L.~J.~Savage, ``Symmetric measures on Cartesian products'',
\emph{Trans.\ Amer.\ Math.\ Soc.}\ 80 (1955), 470--501. P.~L\'evy, \emph{Th\'eorie de l'addition des variables al\'eatoires}, Gauthier-Villars, 1937.
For the story of Borel's paper and its gaps, see G.~Shafer and V.~Vovk, ``The sources of Kolmogorov's Grundbegriffe'', \emph{Statist.\ Sci.}\ 21 (2006), 70--98.
\item R.~Durrett, \emph{Probability: Theory and Examples}, 5th ed., Cambridge University Press, 2019, Sections~2.5 and~4.1. The route followed here for
the zero--one laws and the fate of a random walk.
\item O.~Kallenberg, \emph{Foundations of Modern Probability}, 3rd ed., Springer, 2021, Chapters~4 and~5. Ottaviani's inequality and L\'evy's theorem.
\item J.-P.~Kahane, \emph{Some Random Series of Functions}, 2nd ed., Cambridge University Press, 1985. Random signs, Khintchine's inequalities,
random power series and their natural boundaries.
\item B.~Schmuland, ``Random harmonic series'', \emph{Amer.\ Math.\ Monthly} 110 (2003), 407--416. The density of Figure~\ref{fig:harmonic}, its value
at $0$ and its near miss at $2$.
\end{itemize}
